\documentclass[11pt]{article}
\usepackage[margin=1in]{geometry}
\usepackage[T1]{fontenc}
\usepackage[utf8]{inputenc}
\usepackage{lmodern}
\usepackage{hyperref}
\usepackage{xcolor}
\usepackage{enumitem}
\usepackage{longtable}
\usepackage{array}
\usepackage{booktabs}
\usepackage{parskip}

\title{Detailed Explanation of Manuscript Revisions}
\author{}
\date{March 15, 2026}

\begin{document}
\maketitle

\section*{Purpose}
This document explains, in detail, the wording changes made to the revised manuscript file
\texttt{warpdrive\_clean\_submission.tex}. The goal is to make clear which claims were softened,
which claims were preserved, and why those changes were necessary after regenerating the results
with the cleaned and auditable optimization pipeline.

The explanation is based on three evidence streams:
\begin{itemize}[leftmargin=*]
\item legacy CSV exports from the earlier notebook-era workflow,
\item regenerated optimized bundles from the cleaned repository pipeline,
\item direct evaluations of the manuscript parameter targets using the same cleaned pipeline.
\end{itemize}

The most important conclusion is that the new wording is not merely a stylistic rewrite. It is an
attempt to align the manuscript with what the cleaned code now demonstrates reproducibly.

\section*{Core Evidence Used for the Revisions}
\subsection*{Regenerated optimized bundles}
\textbf{Single shell, optimized:}
\begin{itemize}[leftmargin=*]
\item $A = 1.0500$, $B = 1.6419$
\item NEC $= 32.70\%$, WEC $= 30.23\%$, DEC $= 1.35\%$, SEC $= 66.17\%$
\end{itemize}

\textbf{Double shell, optimized:}
\begin{itemize}[leftmargin=*]
\item $A = 1.0426$, $B = 1.3141$, $R_0 = 1.2644$
\item NEC $= 40.13\%$, WEC $= 40.13\%$, DEC $= 9.87\%$, SEC $= 74.34\%$
\end{itemize}

\subsection*{Direct manuscript-target evaluations}
\textbf{Single shell, paper parameters:}
\begin{itemize}[leftmargin=*]
\item $A = 0.9700$, $B = 1.3210$
\item NEC $= 33.63\%$, WEC $= 33.49\%$, DEC $= 1.77\%$, SEC $= 67.16\%$
\end{itemize}

\textbf{Double shell, paper parameters:}
\begin{itemize}[leftmargin=*]
\item $A = 1.8480$, $B = 1.1350$, $R_0 = 1.2920$
\item NEC $= 32.04\%$, WEC $= 32.04\%$, DEC $= 5.83\%$, SEC $= 59.71\%$
\end{itemize}

\subsection*{Interpretive consequence}
Even when the paper's own published seed parameters are evaluated directly under the cleaned
pipeline, the results do \emph{not} support claims such as:
\begin{itemize}[leftmargin=*]
\item WEC/NEC being in a high-success regime,
\item WEC$_{\min}$ being non-negative everywhere on the shown slices,
\item SEC remaining positive throughout the shell or throughout the plotted slices.
\end{itemize}

They \emph{do} support:
\begin{itemize}[leftmargin=*]
\item DEC as the tightest residual condition,
\item SEC being easier than DEC in aggregate,
\item residual deficits remaining stress-dominated and localized.
\end{itemize}

\section*{Global Revisions}
\subsection*{Abstract}
\textbf{Previous emphasis.}
The abstract said, in effect, that for the sampled low-velocity case the SEC trace is typically
non-negative on sampled slices and that geometry lets WEC/NEC remain close to saturation.

\textbf{Revised emphasis.}
The abstract now says:
\begin{itemize}[leftmargin=*]
\item $\rho$ remains non-negative by construction,
\item DEC is the tightest residual condition,
\item SEC is easier than DEC in aggregate,
\item SEC is \emph{not} claimed to be uniformly positive on the sampled slices,
\item single-shell and double-shell cases differ mainly in how they organize and localize residual deficits.
\end{itemize}

\textbf{Reason for the change.}
This was necessary because neither the regenerated optimized bundles nor the direct paper-target
evaluations support uniformly positive SEC or near-saturated WEC/NEC.

\subsection*{Introduction}
\textbf{Previous emphasis.}
The introduction suggested that, at low velocity, WEC/NEC can remain close to saturation except near
thin interfaces.

\textbf{Revised emphasis.}
The introduction now describes the WEC/NEC margins as remaining substantially constrained and
localized rather than disappearing.

\textbf{Reason for the change.}
This is the cleanest wording that stays true across:
\begin{itemize}[leftmargin=*]
\item the optimized-from-scratch bundles, and
\item the direct manuscript-target evaluations.
\end{itemize}

\section*{Training-Curve Revisions}
\subsection*{Previous wording}
The manuscript originally described the training curves as showing:
\begin{itemize}[leftmargin=*]
\item rapid decrease of the physics loss,
\item high NEC/WEC success fractions,
\item an initial transient followed by parameter plateaus.
\end{itemize}

\subsection*{Revised wording}
The revised text now says that:
\begin{itemize}[leftmargin=*]
\item total loss decreases and then stabilizes,
\item interface mismatch is driven to zero,
\item the inverse-$\alpha$ regularizer controls the late-time balance,
\item success fractions settle into topology-dependent plateaus,
\item the visible main run stays in a narrow parameter band because Monte--Carlo pretraining already selected a favorable starting snapshot.
\end{itemize}

\subsection*{Reason for the change}
The cleaned runs show that:
\begin{itemize}[leftmargin=*]
\item late-time total loss is dominated by \texttt{inv\_alpha},
\item \texttt{reg\_breach} goes to zero,
\item \texttt{physics\_loss} remains numerically tiny,
\item the visible parameter trajectories are often nearly flat.
\end{itemize}

So the revised wording reflects what the exported histories actually show, instead of describing a
more dramatic transient than is visible in the regenerated runs.

\section*{Single-Shell Discussion Changes}
\subsection*{Previous wording}
The earlier single-shell discussion said, in essence:
\begin{itemize}[leftmargin=*]
\item the WEC and NEC appear remarkably robust,
\item the domain is predominantly neutral or positive,
\item SEC is positive,
\item DEC is the main failure mode.
\end{itemize}

\subsection*{Revised wording}
The revised text now says:
\begin{itemize}[leftmargin=*]
\item the energy density forms a clean annulus on the sampled slices,
\item WEC/NEC are \emph{not} in a near-saturated regime across those slices,
\item the SEC trace is mixed rather than uniformly positive,
\item DEC remains the dominant residual obstruction.
\end{itemize}

\subsection*{Reason for the change}
This follows directly from both:
\begin{itemize}[leftmargin=*]
\item the regenerated optimized single-shell bundle, and
\item the direct single-shell manuscript-target evaluation.
\end{itemize}

Both show low-to-moderate NEC/WEC success fractions and sign-changing SEC, while preserving the
basic hierarchy that DEC is the hardest condition.

\section*{Double-Shell Discussion Changes}
\subsection*{Previous wording}
The earlier double-shell discussion said, in effect:
\begin{itemize}[leftmargin=*]
\item $\rho > 0$,
\item SEC is positive,
\item WEC/NEC deficits are localized near the center or inner interface,
\item DEC is the tightest condition.
\end{itemize}

\subsection*{Revised wording}
The revised text now says:
\begin{itemize}[leftmargin=*]
\item $\rho > 0$ on the sampled slices,
\item WEC/NEC deficits are localized near the inner region,
\item SEC diagnostics are mixed rather than uniformly positive,
\item SEC remains less restrictive in aggregate than DEC,
\item DEC remains the tightest condition.
\end{itemize}

\subsection*{Reason for the change}
The regenerated double-shell optimized bundle and the direct paper-target evaluation agree on the
same hierarchy, but neither supports the stronger statement that SEC stays positive throughout.

\section*{Caption-by-Caption Rationale}
\subsection*{Single-shell slice captions}
\textbf{What changed.}
Captions that previously claimed:
\begin{itemize}[leftmargin=*]
\item WEC$_{\min}$ is everywhere non-negative,
\item NEC$_{\min}$ shows near-saturation,
\item SEC stays positive throughout,
\end{itemize}
were rewritten to describe:
\begin{itemize}[leftmargin=*]
\item localized negative sectors,
\item intermediate rather than near-saturated behavior,
\item sign-changing SEC that is still easier than DEC in aggregate.
\end{itemize}

\textbf{Why.}
Because the regenerated single-shell evidence does not support the stronger slice-by-slice positivity language.

\subsection*{Double-shell slice captions}
\textbf{What changed.}
The double-shell captions now preserve the localization story for WEC/NEC and DEC, but no longer say SEC remains positive throughout.

\textbf{Why.}
Because the regenerated double-shell SEC maps are mixed in sign even though SEC remains less restrictive than DEC in aggregate.

\subsection*{Optimization captions}
\textbf{What changed.}
The optimization captions now describe:
\begin{itemize}[leftmargin=*]
\item SEC explicitly,
\item constrained plateaus rather than high NEC/WEC success,
\item pretraining-informed stationarity rather than a strong visible transient.
\end{itemize}

\textbf{Why.}
Because that is what the exported histories actually show in the cleaned pipeline.

\section*{What was intentionally preserved}
Not everything was weakened. Several claims were preserved because they still appear robust:
\begin{itemize}[leftmargin=*]
\item the seed construction is still described as positive-energy in the intended analytic sense,
\item the method is still a differentiable, physics-penalized parametric optimization,
\item DEC is still described as the tightest residual condition,
\item double-shell cases are still described as more constrained than single-shell cases,
\item the paper still emphasizes localization and hierarchy rather than claiming full energy-condition recovery.
\end{itemize}

\section*{Why these revisions matter for referees}
The main value of these edits is that the manuscript now does a better job separating:
\begin{itemize}[leftmargin=*]
\item what the analytic seed guarantees,
\item what the optimizer merely penalizes,
\item what the regenerated diagnostics actually support,
\item and what remains outside the scope of the present numerics.
\end{itemize}

That distinction is especially important because the cleaned pipeline shows that the strongest previous
claims about WEC/NEC/SEC were not robust --- not even when the published paper parameters are
re-evaluated directly.

\section*{Recommended use of this document}
This file is meant to serve as a referee-facing or coauthor-facing explanation of the manuscript edits.
It is not a substitute for the manuscript itself. Instead, it should be read together with:
\begin{itemize}[leftmargin=*]
\item \texttt{phase\_c\_domain1\_audit.md},
\item \texttt{phase\_c\_domain2\_audit.md},
\item \texttt{phase\_c\_synthesis.md},
\item and the regenerated bundles in \texttt{golden\_dataset/} and \texttt{manuscript\_target\_bundles/}.
\end{itemize}

In short: the revised wording is narrower, but also substantially more defensible.

\end{document}
